Music has accompanied human language for as long as recorded history, and the
song --- lyrics set to melody --- remains one of the most universal forms of
creative expression. Composing a melody for a given lyric, however, is a task
that has traditionally demanded years of musical training. It requires the
composer to understand what the words \emph{mean}, to feel what they
\emph{evoke}, and to translate both into a sequence of pitches and durations
that a human voice can sing naturally. This project investigates whether a
Large Language Model (LLM), used purely through prompting and without any
task-specific training, can perform this act of cross-modal composition, and
presents \textbf{L2M}, a complete, open-source lyrics-to-melody generation
system built on this premise.

\section{Motivation}
\label{sec:motivation}

Automatic music composition has advanced rapidly with deep learning, and
systems now exist that generate convincing instrumental music
\cite{briot20deeplearning, huang18musictransformer} and even high-fidelity
audio directly from text descriptions \cite{agostinelli23musiclm,
copet23musicgen}. Yet the specific task of \emph{lyrics-to-melody generation}
--- producing a singable melodic line for a given lyrical text --- remains
comparatively difficult, because it couples natural language understanding
with symbolic music generation under hard structural constraints.

Existing lyrics-to-melody systems \cite{bao19neural, sheng21songmass,
ju22telemelody, zhang22relyme, ding24songcomposer} are almost exclusively
supervised: they learn the mapping from lyrics to melody from large corpora
of paired examples. Such corpora are scarce, expensive to curate, and often
encumbered by copyright, which restricts both research reproducibility and
practical deployment. Moreover, purely neural systems typically offer no
recourse when the model produces an invalid output at inference time ---
a wrong number of notes, an out-of-range pitch, or a malformed sequence
simply becomes a failed run.

The emergence of instruction-following LLMs trained on internet-scale corpora
--- which include music theory texts, song analyses, and transcribed lyrics
--- suggests an alternative path. If an LLM already encodes substantial
implicit knowledge about the relationship between words, emotions, and music,
then careful \emph{prompt engineering} may elicit that knowledge zero-shot
\cite{brown20fewshot}, eliminating the need for paired training data
altogether. Whether this is actually achievable, and how reliably, is the
central question motivating this work.

\section{Problem Statement}
\label{sec:problem-statement}

The lyrics-to-melody task imposes four tightly coupled requirements on any
generation system:

\begin{enumerate}
    \item \textbf{Semantic understanding.} The system must interpret not
    merely the denotative meaning of the words but the affective intent and
    thematic atmosphere embedded in the text.
    \item \textbf{Syllabic alignment.} Every phonetic syllable must
    correspond to exactly one musical note, preserving the natural cadence
    of the sung phrase and ensuring singability.
    \item \textbf{Musical coherence.} The output must exhibit a well-chosen
    key, consistent tempo, appropriate mode, and a melodic contour that forms
    a satisfying phrase arc within the conventions of tonal music.
    \item \textbf{Emotional consistency.} The affective character of the
    melody must reinforce, rather than contradict, the sentiment expressed
    in the lyrics.
\end{enumerate}

Formally, given an arbitrary lyrical text $L$ containing $N$ syllables, the
system must produce a melody $M = \langle (p_1, d_1), \ldots, (p_N, d_N)
\rangle$ --- a sequence of exactly $N$ (pitch, duration) pairs --- together
with a key, tempo, and time signature, such that $M$ is singable, tonally
coherent, and emotionally consistent with $L$. The system must produce a
valid $M$ for \emph{every} input, including when the underlying LLM service
is unavailable or returns malformed output.

\section{Research Questions}
\label{sec:research-questions}

This project addresses the following research questions:

\begin{itemize}
    \item[\textbf{RQ1.}] Can a general-purpose LLM, accessed zero-shot through
    structured prompts, generate melodies that satisfy exact syllable--note
    alignment and basic tonal coherence, without any task-specific training?
    \item[\textbf{RQ2.}] Does decomposing the task into two stages ---
    emotion/rhythm analysis followed by conditioned melody generation ---
    improve controllability and interpretability over a single monolithic
    generation step?
    \item[\textbf{RQ3.}] Can a deterministic, music-theory-grounded fallback
    mechanism guarantee full operational reliability while preserving
    acceptable musical quality when the LLM path fails?
    \item[\textbf{RQ4.}] Which prompt engineering techniques (output schemas,
    hard constraints, few-shot examples, contextual carryover) are necessary
    to make LLM output usable in a strictly validated music pipeline?
\end{itemize}

\section{Objectives}
\label{sec:objectives}

The concrete objectives of this work are to:

\begin{enumerate}
    \item Design an end-to-end pipeline that accepts arbitrary lyrical text
    and produces standard music notation (MIDI, MusicXML) and optional
    synthesized audio (WAV/MP3);
    \item Develop a two-stage LLM prompting strategy for emotion analysis
    and syllable-aligned melody generation, with structured JSON output
    schemas validated at every stage;
    \item Ground the generation process in music theory through an explicit
    mapping from emotion categories to musical keys, tempo ranges, and
    melodic contour patterns;
    \item Implement deterministic fallback heuristics that guarantee a valid
    output whenever an LLM call fails, times out, or returns malformed data;
    \item Evaluate the system quantitatively and qualitatively on a benchmark
    spanning six emotion categories, and release it as a documented,
    open-source Python package.
\end{enumerate}

\section{Contributions}
\label{sec:contributions}

The principal contributions of this project are:

\begin{enumerate}
    \item \textbf{Two-stage LLM pipeline.} A novel decomposition of the
    lyrics-to-melody task into (a) emotion and rhythm analysis and (b)
    conditioned melody generation, each handled by a separate LLM call with
    a structured JSON output schema --- enabling interpretability and
    targeted fallback at each stage.
    \item \textbf{Emotion-aware alignment.} An explicit mapping from detected
    emotion categories to musical keys, tempo ranges, and melodic contour
    patterns, grounding LLM-generated output in established music-psychology
    findings \cite{hevner36expression, juslin01musicemotion}.
    \item \textbf{Hybrid reliability architecture.} A combination of
    zero-shot LLM generation with deterministic heuristic fallbacks that
    achieves a 100\% completion rate across all tested inputs.
    \item \textbf{Prompt engineering methodology.} Structured few-shot
    prompts with explicit output schemas, hard constraint enforcement
    (``EXACTLY $N$ notes for $N$ syllables''), and contextual carryover
    across long-lyric chunks --- a reusable template for other cross-modal
    LLM tasks.
    \item \textbf{Open-source baseline.} A production-ready Python package
    offering a CLI, a programmatic API, multi-format output,
    and comprehensive logging, providing a reproducible baseline for the
    research community.
\end{enumerate}

\section{Scope and Limitations}
\label{sec:scope}

The system generates \emph{monophonic} vocal melodies within the conventions
of Western tonal music; harmony, accompaniment, and non-Western musical
systems are outside the present scope and are discussed as future work.
Evaluation is conducted on a 20-sample proof-of-concept benchmark with
automatic metrics and expert qualitative review; a large-scale listener
study is likewise deferred to future work.

\section{Report Organization}
\label{sec:organization}

The remainder of this report is organized as follows.
\Autoref{chap:literature-review} surveys prior work on lyrics-to-melody
generation, text-to-music synthesis, and LLM-based generation.
\Autoref{chap:background} introduces the background concepts this work
builds upon: transformer-based LLMs, prompt engineering, the psychology of
musical emotion, and symbolic music representations.
\Autoref{chap:methodology} presents the proposed six-stage L2M architecture
in detail, including both LLM prompt designs and the deterministic fallback
algorithms. \Autoref{chap:results-analysis} reports the experimental setup,
quantitative results, qualitative analyses, and a comparison with related
systems, followed by a discussion of findings and limitations.
\Autoref{chap:conclusion-future-work} concludes and outlines directions for
future research.
